\chapter*{\large ЗАКЛЮЧЕНИЕ}
\addcontentsline{toc}{chapter}{ЗАКЛЮЧЕНИЕ}

В ходе курсовой работы была создана программная система для игры <<Змейка>> и реализован агент, который обучается методом Q-learning. В работе были разобраны основы обучения с подкреплением, спроектирована игровая среда, выбрано описание состояния и задана система наград.

Практическая часть показала, что даже табличный подход позволяет агенту постепенно улучшать свою игру. В вычислительных экспериментах средний результат после 100 эпизодов обучения составил 7.544 очка, после 500 эпизодов --- 13.958 очка, а после 1200 эпизодов --- 16.230 очка. Средний лучший результат серии вырос от 15.800 до 21.600 очка. Это показывает, что по мере обучения агент действительно начинает принимать более удачные решения, хотя на поздних этапах рост качества постепенно замедляется.

Можно сделать вывод, что поставленная цель достигнута: игра создана, обучаемый агент реализован, а его работа проанализирована. Полученные результаты показывают, что выбранное описание состояния, относительные действия и система наград позволяют получить полезную стратегию даже без нейронных сетей.

С точки зрения машинного обучения важно, что успех агента связан не только с самим правилом обновления Q-learning, но и с тем, как построена среда. В проекте используется вектор из 11 бинарных признаков, который хранит ключевую информацию о ближайшей опасности, текущем направлении и положении пищи. За счёт этого пространство состояний остаётся компактным для табличного метода.

Большую роль играет и система наград. Положительная награда за пищу, штраф за столкновение и небольшой штраф за обычный шаг помогают агенту учиться не просто выживать, а искать полезные траектории. Также важен баланс между исследованием и использованием уже найденных удачных действий: сначала агент чаще пробует разные варианты, а затем постепенно закрепляет более сильную стратегию.

Дальше работу можно развивать в нескольких направлениях. Можно подробнее исследовать влияние параметров $\alpha$, $\gamma$ и $\varepsilon$, попробовать другое описание состояния, увеличить размер поля или перейти к более сложным методам, например DQN. Такой переход особенно интересен в тех случаях, когда табличный Q-learning уже плохо справляется из-за слишком большого пространства состояний.
